%% Part V -- Risks (comic pages 37 to 43).
%% Five risks, and none of them is the take-over the Intro opened with. They
%% are surveillance, a fragmented public sphere, the labour market, data
%% security, and the energy bill. Each is already running.
%% Fragment only: no preamble, \input by ai_for_all.tex.

\sectionsubtitle{pages 37 to 43}
\section{Risks}

% Page 39, the through-wall radio work. The page walks from words and faces
% to gait to radio reflections, and the step that matters is the last one:
% identity without line of sight.
\begin{frame}{Surveillance}
  \panel{page 39, the stick figures behind the wall}
  \begin{comicquote}{39}
    That means the AI does not only know where we are. The AI can also
    distinguish what we are doing, whether we are standing, or eating.
  \end{comicquote}
  \slidesources{Schneider2019}
\end{frame}

% Page 40. The risk named is structural, not a bad actor: the public sphere
% has broken into pieces, and the pieces are privately owned.
\begin{frame}{Opinion-Forming and Media}
  \panel{page 40, the fragmented public sphere}
  \begin{comicquote}{40}
    AI influences public opinion-forming and the media. Our public sphere
    has become fragmented.
  \end{comicquote}
  \slidesources{Schneider2019}
\end{frame}

% Page 41 is the labour page, and it refuses the usual blue-collar framing:
% lawyers, medics and brokers are on the list next to drivers. The quote
% taken is the other half, what stays with us and why.
\begin{frame}{Future Work}
  \panel{page 41, the job ladder}
  \begin{comicquote}{41}
    We are better in understanding non-trivial, new situations, building
    relationships, creating context and meaning. Custom-fit solutions.
  \end{comicquote}
  \slidesources{Schneider2019}
\end{frame}

% Page 42 is the CCC programme: the data letter, penalties for outflows,
% declared update frequency, permanent state investment in secure open
% source, and education. The opening sentence is the one that earns them.
\begin{frame}{Data Security and Safety}
  \panel{page 42, the data letter inbox}
  \begin{comicquote}{42}
    Our blind trust in AI providers is misplaced. We all know that the
    amount of data that we are implicitly or explicitly forced to leave
    behind is too great.
  \end{comicquote}
  \slidesources{Schneider2019}
\end{frame}

% Page 43. The comparison is with a brain, not with a previous data centre,
% and the numbers on the page (33 zettabyte in 2018, 175 by 2025, 8\% of
% global energy by 2030) are the ones to read out next to the quote.
\begin{frame}{Energy Use}
  \panel{page 43, save energy, save the planet}
  \begin{comicquote}{43}
    But conventional hardware was not designed for AI, hence it needs much
    more computing power, much more energy, than our brains do.
  \end{comicquote}
  \slidesources{Schneider2019}
\end{frame}
